\chapter{Resultados y evaluación}\label{cap:resultados_y_evaluacion}

Una vez finalizado el ciclo de desarrollo e implementadas las pruebas de verificación expuestas en el capítulo anterior, el presente capítulo aborda la auditoría cuantitativa y cualitativa del proyecto. A través de este análisis se evalúa el desempeño de la solución construida respecto a los compromisos definidos durante la fase de planificación, contrastando la ejecución real frente a las líneas base establecidas. Asimismo, se documenta la gestión de cambios articulada durante el desarrollo, la resolución efectiva de las pruebas y una retrospectiva técnica en torno a la infraestructura de integración continua y métricas de calidad de código.
 
\section{Cumplimiento de la Línea Base del Alcance}\label{sec:cumplimiento_alcance}

La \hyperref[sec:linea_base_del_alcance]{Línea Base del Alcance} quedó formalizada en la \hyperref[cap:gestion_del_proyecto]{planificación del proyecto} a través de la \hyperref[subsec:edt]{EDT} y su correspondiente \hyperref[subsec:diccionario_edt]{diccionario}, distribuyendo el esfuerzo en cuatro paquetes troncales: Gestión del proyecto, Diseño y arquitectura, Producto software y Documentación final. 

A nivel global del proyecto, los entregables documentales, de gestión y de diseño preliminar se han completado al \textbf{100\%}, quedando consolidados a lo largo de los capítulos técnicos, anexos y código fuente de esta memoria. Por su parte, el paquete \textbf{Producto software} materializa con éxito el flujo íntegro del sistema a través de las tres \textit{releases} planificadas; no obstante, al haberse pospuesto dos de los doce objetivos funcionales contemplados en el diseño inicial, el grado de consecución de dichos objetivos se sitúa en el \textbf{83{,}33\%}.

Para auditar las desviaciones del alcance del sistema, el siguiente cuadro detalla en primer término los dos objetivos funcionales que resultaron formalmente pospuestos, justificando el motivo de su aplazamiento:

\begin{table}[H]
\centering
\normalsize
\renewcommand{\arraystretch}{1.5}
\makebox[\textwidth][c]{%
\begin{tabular}{llp{10.7cm}}
\hline
\textbf{ID} & \textbf{Objetivo funcional} & \textbf{Motivo del aplazamiento} \\ \hline
\hyperref[obj:09]{OBJ-09} & Asistente IA & Evitar la dependencia de cuotas de facturación por consumo de \textit{tokens}, límites de tasa (\textit{rate limits}) y sobrecostes imprevistos de servicios externos, priorizando un núcleo completamente autosuficiente. \\ 
\hyperref[obj:11]{OBJ-11} & Gestión de grupos & Elevado coste de oportunidad: el gran esfuerzo temporal necesario para su implementación no justificaba el valor añadido, decidiendo concentrar ese tiempo en consolidar la concurrencia y la estabilidad de las salas en vivo. \\ \hline
\end{tabular}%
}
\caption{Objetivos funcionales pospuestos.}
\label{tab:objetivos_pospuestos}
\end{table}

Gracias a las relaciones establecidas en la \hyperref[tab:matriz_trazabilidad]{matriz de trazabilidad unificada}, es posible determinar con precisión cómo repercute la postergación de estos dos objetivos en los casos de uso, historias de usuario y requisitos funcionales asociados:

\begin{itemize}
    \item \textbf{Asistente de IA (\hyperref[obj:09]{OBJ-09}):} Conlleva la postergación completa del caso de uso \textbf{\hyperref[cu:03]{CU-03} - Generar preguntas con IA} y de la historia de usuario \textbf{\hyperref[hu:pr-09]{HU-PR-09} - Generación asistida de preguntas}, afectando a cuatro requisitos funcionales. Estos requisitos han quedado formalizados en el repositorio del proyecto como líneas de evolución futura:

\begin{table}[H]
\centering
\normalsize
\renewcommand{\arraystretch}{1.5}

\makebox[\textwidth][c]{%
\begin{tabular}{llc}
\hline
\textbf{ID} & \textbf{Requisito funcional} & \textbf{GitHub Issue} \\ \hline
\hyperref[req:34]{RF-34} & Generación preguntas & \issue{95}  \\
\hyperref[req:35]{RF-35} & Asistente de ayuda &  \issue{96} \\
\hyperref[req:36]{RF-36} & Navegación y tutorial por IA & \issue{97} \\
\hyperref[req:37]{RF-37} & Revisión post IA & \issue{98} \\ \hline
\end{tabular}%
}
\caption{Requisitos funcionales pospuestos para OBJ-09.}
\label{tab:rf_pospuestos_obj09}
\end{table}

    \item \textbf{Gestión de grupos (\hyperref[obj:11]{OBJ-11}):} Implica posponer íntegramente el caso de uso \textbf{\hyperref[cu:11]{CU-11} - Gestionar grupos} y la historia de usuario \textbf{\hyperref[hu:pr-07]{HU-PR-07} - Organización por grupos de clase}, afectando a un único requisito funcional:
\begin{table}[H]
\centering
\normalsize
\renewcommand{\arraystretch}{1.5}
\makebox[\textwidth][c]{%
\begin{tabular}{llc}
\hline
\textbf{ID} & \textbf{Requisito funcional} & \textbf{GitHub Issue} \\ \hline
\hyperref[req:38]{RF-38} & Crear clase & \issue{99} \\ \hline
\end{tabular}%
}
\caption{Requisito funcional pospuesto para OBJ-11.}
\label{tab:rf_pospuestos_grupos}
\end{table}

\end{itemize}

Por otro lado, los diez objetivos funcionales restantes alcanzaron cobertura operativa suficiente para la entrega. La siguiente tabla resume la retrospectiva técnica, los desafíos de diseño y las decisiones clave adoptadas durante su materialización:

\begin{table}[H]
\centering
\normalsize
\renewcommand{\arraystretch}{1.5}
\makebox[\textwidth][c]{%
\begin{tabular}{lp{3.7cm}p{10.7cm}}
\hline
\textbf{ID} & \textbf{Objetivo funcional} & \textbf{Retrospectiva técnica e impacto en el desarrollo} \\ \hline
\hyperref[obj:01]{OBJ-01} & Módulo de usuario y autenticación & Implementación sólida de JWT y Bcrypt. Para evitar bloqueos en el hilo principal de la API por la latencia del servidor SMTP de Gmail, el envío del código de verificación por correo se delegó a tareas en segundo plano. \\ 
\hyperref[obj:02]{OBJ-02} & Gestión de cuestionarios & Se optó por borrado lógico (\textit{soft delete}) en cascada para impedir que la edición o eliminación de preguntas corrompiera los registros de salas históricas. \\ 
\hyperref[obj:03]{OBJ-03} & Creación de salas & Se implementó una máquina de estados en el servidor para controlar las fases de la partida, garantizando un flujo sincronizado y previniendo saltos de estado indebidos. \\ 
\hyperref[obj:04]{OBJ-04} & Acceso rápido & Entrada ágil e inmediata sin necesidad de registro previo. Para evitar que un refresco accidental del navegador expulsara al estudiante de la partida, la sesión activa se conservó localmente mediante \texttt{sessionStorage}. \\ 
\hyperref[obj:05]{OBJ-05} & Interacción en tiempo real & La recepción de respuestas se gestionó en memoria garantizando que cada participante emitiera un único voto por pregunta. Además, se centralizó el temporizador en el servidor para anular intentos de manipulación o ventajas por desajustes horarios entre los clientes. \\ 
\hyperref[obj:06]{OBJ-06} & Verificación QR & Generación distribuida en el cliente para liberar de procesamiento al servidor mediante la librería \texttt{qrcode.react}. La carga útil del código QR estructura un objeto JSON con los datos clave de la evaluación (UVUS, sala y puntuación), acompañado de un token criptográfico de un solo uso. \\ 
\hyperref[obj:07]{OBJ-07} & Persistencia validada & Se logró el desacoplamiento entre el estado efímero de la partida y el almacenamiento persistente: solo los resultados verificados se consolidan en la base de datos. \\  
\hyperref[obj:08]{OBJ-08} & Estadísticas y análisis & Agregación reactiva en memoria al cierre de cada ronda para computar la distribución de respuestas, aciertos globales y el podio en tiempo real, evitando sobrecargar la base de datos con consultas recurrentes. \\  
\hyperref[obj:10]{OBJ-10} & Escaneo de código QR & Lectura ágil mediante \texttt{html5-qrcode} con solicitud de permisos de cámara, validando el \textit{token} criptográfico contra el servidor en un solo paso atómico. \\  
\hyperref[obj:12]{OBJ-12} & Exportación de datos & Se habilitó la descarga de resultados en formato CSV, garantizando la compatibilidad con hojas de cálculo. \\  \hline
\end{tabular}%
}
\caption{Retrospectiva técnica de los objetivos funcionales.}
\label{tab:objetivos_completados}
\end{table}

\newpage

Cabe destacar que los nueve casos de uso y dieciséis historias de usuario que sustentan estos diez objetivos están plenamente disponibles en sus flujos troncales. Sin embargo, para cumplir con el calendario fijado sin comprometer la estabilidad, se decidió catalogar como mejoras o extensiones secundarias siete requisitos funcionales. La siguiente tabla detalla estos requisitos de ampliación y las \textit{issues} abiertas para su integración en versiones sucesivas:

\begin{table}[H]
\centering
\normalsize
\renewcommand{\arraystretch}{1.5}
\makebox[\textwidth][c]{%
\begin{tabular}{llc}
\hline
\textbf{ID} & \textbf{Requisito funcional} & \textbf{GitHub Issue} \\ \hline
\hyperref[req:12]{RF-12} & Importación & \issue{100} \\
\hyperref[req:32]{RF-32} & Listado en vivo & \issue{101} \\
\hyperref[req:33]{RF-33} & Estadísticas & \issue{102} \\
\hyperref[req:42]{RF-42} & Bonificación por tiempo & \issue{103} \\
\hyperref[req:43]{RF-43} & Sistema de rachas & \issue{104} \\
\hyperref[req:44]{RF-44} & Configuración de puntuación & \issue{105} \\
\hyperref[req:46]{RF-46} & Creación/edición avanzada & \issue{106} \\ \hline
\end{tabular}%
}
\caption{Requisitos funcionales pospuestos como puntos de extensión.}
\label{tab:rf_mejoras_issues}
\end{table}

En términos cuantitativos, este ajuste de alcance deliberado se traduce en que 9 de los 11 casos de uso diseñados (\textbf{81{,}82\,\%}) y 16 de las 18 historias de usuario (\textbf{88{,}89\,\%}) se encuentran plenamente operativos en sus flujos de valor principal. En el nivel de especificación atómica, 35 de los 47 requisitos funcionales (\textbf{74{,}47\,\%}) han sido implementados y verificados con éxito, garantizando la totalidad de la operativa requerida en el aula presencial y dejando acotadas las extensiones como trabajo futuro.

\section{Tiempo real frente a Línea Base del Cronograma}\label{sec:tiempo_real_vs_cronograma}

Para auditar el desempeño temporal del proyecto frente a la \hyperref[sec:linea_base_del_cronograma]{Línea Base del Cronograma}, se ha consolidado el registro continuo de dedicación temporal a través de Clockify, cuyo informe detallado queda recogido íntegramente en el \hyperref[anexo:clockify]{anexo de dedicación temporal}. Las horas se han computado a nivel de tareas operativas individuales, asignando en la herramienta una etiqueta (\textit{tag}) que las vincula directamente con su actividad homóloga más próxima de la planificación inicial. Este mecanismo permite clasificar el esfuerzo según el rol profesional desempeñado, trazabilidad que se empleará más adelante en la \hyperref[tab:horas_roles_previstas_reales]{tabla de dedicación por roles} para calcular las desviaciones en los gastos de contratación.

\textcolor{red}{\textbf{[TODO: Actualizar informe cuando estén todas las horas computadas]}}

El siguiente cuadro resume la comparativa entre la estimación inicial y la dedicación real registrada por fase de trabajo:

\textcolor{red}{\textbf{[TODO: Rellenar las horas]}}

\begin{table}[H]
\centering
\small
\renewcommand{\arraystretch}{1.3}
\begin{tabular}{llcc}
\hline
\textbf{ID} & \textbf{Fase} & \textbf{Horas previstas} & \textbf{Horas reales} \\ \hline
F-01 & Planificación & 20 h & 11 h \\
F-02 & Documentación & 20 h & 27,14 h \\
F-03 & Inicio        & 11 h & 12,50 h \\
F-04 & CI/CD         & 17 h & 19,87 h \\
F-05 & MVP           & 65 h & 55,40 h \\
F-06 & Testing       & 15 h & 42,74 h \\
F-07 & Core          & 49 h & 45,24 h \\
F-08 & QA            & 13 h & 23,73 h \\
F-09 & Memoria       & 70 h &   \\
F-10 & Defensa       & 20 h &  \\ \hline
\multicolumn{2}{l}{\textbf{Total}} & \textbf{300 h} &   \\ \hline
\end{tabular}
\caption{Comparativa de horas previstas y reales por fase del proyecto.}
\label{tab:horas_fases_previstas_reales}
\end{table}

\textcolor{red}{\textbf{[TODO: Redactar retrospectiva general de desviaciones temporales una vez cerrado el cómputo final de horas en Clockify]}}

Este incremento en el esfuerzo dedicado a la robustez técnica y las pruebas del sistema impactó de forma directa en los plazos globales. Si bien la planificación inicial contemplaba finalizar el proyecto para la convocatoria ordinaria de junio, mantener dicha fecha habría obligado a comprometer la calidad del producto o recortar el rigor de la memoria técnica. Por ello, se optó formalmente por trasladar la finalización y defensa a la convocatoria extraordinaria de octubre, quedando esta modificación justificada y registrada en la sección de \hyperref[sec:desviaciones_y_gestion_cambios]{gestión de cambios}.

A continuación, se muestra una comparativa entre la Línea Base del Cronograma y la ejecución real, representada mediante un diagrama de Gantt:

\begin{figure}[H]
\centering
\noindent\makebox[\textwidth][c]{%
\begin{ganttchart}[
    hgrid=false, 
    vgrid=false, 
    x unit=0.0636cm,
    y unit title=0.55cm,
    y unit chart=0.72cm, 
    time slot format=isodate,
    inline,
    title label font=\scriptsize\bfseries,
    title calendar link/.style={draw=black},
    canvas/.style={draw=black, fill=white}
]{2026-01-15}{2026-10-17}

    \gantttitle{Ene}{17}
    \gantttitle{Feb}{28}
    \gantttitle{Mar}{31}
    \gantttitle{Abr}{30}
    \gantttitle{May}{31}
    \gantttitle{Jun}{30}
    \gantttitle{Jul}{31}
    \gantttitle{Ago}{31}
    \gantttitle{Sep}{30}
    \gantttitle{Oct}{17} \\

    % F01 
    \ganttbar[bar/.style={fill=gray!45, draw=primaryMagenta, line width=0.5pt, rounded corners=1pt}, bar height=0.40, bar top shift=0.08, bar label font=\tiny\bfseries\color{black!80}, ganttbar label node/.append style={xshift=0.35cm}]{F01-Plan}{2026-01-15}{2026-01-28}
    \ganttbar[bar/.style={fill=primaryMagenta, draw=primaryMagenta, rounded corners=1pt}, bar height=0.40, bar top shift=0.48, bar label font=\tiny\bfseries\color{white}, ganttbar label node/.append style={xshift=0.15cm}]{F01}{2026-01-15}{2026-01-26} \\

    % F02  
    \ganttbar[bar/.style={fill=gray!45, draw=secondaryMagenta, line width=0.5pt, rounded corners=1pt}, bar height=0.40, bar top shift=0.08, bar label font=\tiny\bfseries\color{black!80}]{F02-Plan}{2026-01-29}{2026-04-14}
    \ganttbar[bar/.style={fill=secondaryMagenta, draw=secondaryMagenta, rounded corners=1pt}, bar height=0.40, bar top shift=0.48, bar label font=\tiny\bfseries\color{white}]{F02}{2026-01-26}{2026-06-07} \\

    % F03 
    \ganttbar[bar/.style={fill=gray!45, draw=primaryMagenta, line width=0.5pt, rounded corners=1pt}, bar height=0.40, bar top shift=0.08, bar label font=\tiny\bfseries\color{black!80}]{F03-Plan}{2026-01-29}{2026-02-11}
    \ganttbar[bar/.style={fill=primaryMagenta, draw=primaryMagenta, rounded corners=1pt}, bar height=0.40, bar top shift=0.48, bar label font=\tiny\bfseries\color{white}]{F03}{2026-01-30}{2026-02-09} \\

    % F04  
    \ganttbar[bar/.style={fill=gray!45, draw=secondaryCyan, line width=0.5pt, rounded corners=1pt}, bar height=0.40, bar top shift=0.08, bar label font=\tiny\bfseries\color{black!80}]{F04-Plan}{2026-02-12}{2026-04-26}
    \ganttbar[bar/.style={fill=secondaryCyan, draw=secondaryCyan, rounded corners=1pt}, bar height=0.40, bar top shift=0.48, bar label font=\tiny\bfseries\color{white}]{F04}{2026-02-11}{2026-08-14} \\

    % F05  
    \ganttbar[bar/.style={fill=gray!45, draw=primaryCyan, line width=0.5pt, rounded corners=1pt}, bar height=0.40, bar top shift=0.08, bar label font=\tiny\bfseries\color{black!80}]{F05-Plan}{2026-02-19}{2026-03-08}
    \ganttbar[bar/.style={fill=primaryCyan, draw=primaryCyan, rounded corners=1pt}, bar height=0.40, bar top shift=0.48, bar label font=\tiny\bfseries\color{white}]{F05}{2026-02-12}{2026-04-04} \\

    % F06 
    \ganttbar[bar/.style={fill=gray!45, draw=secondaryCyan, line width=0.5pt, rounded corners=1pt}, bar height=0.40, bar top shift=0.08, bar label font=\tiny\bfseries\color{black!80}]{F06-Plan}{2026-02-24}{2026-04-26}
    \ganttbar[bar/.style={fill=secondaryCyan, draw=secondaryCyan, rounded corners=1pt}, bar height=0.40, bar top shift=0.48, bar label font=\tiny\bfseries\color{white}]{F06}{2026-04-06}{2026-08-14} \\

    % F07 
    \ganttbar[bar/.style={fill=gray!45, draw=primaryCyan, line width=0.5pt, rounded corners=1pt}, bar height=0.40, bar top shift=0.08, bar label font=\tiny\bfseries\color{black!80}]{F07-Plan}{2026-03-09}{2026-04-12}
    \ganttbar[bar/.style={fill=primaryCyan, draw=primaryCyan, rounded corners=1pt}, bar height=0.40, bar top shift=0.48, bar label font=\tiny\bfseries\color{white}]{F07}{2026-04-15}{2026-04-28} \\

    % F08 
    \ganttbar[bar/.style={fill=gray!45, draw=primaryCyan, line width=0.5pt, rounded corners=1pt}, bar height=0.40, bar top shift=0.08, bar label font=\tiny\bfseries\color{black!80}]{F08-Plan}{2026-04-13}{2026-04-24}
    \ganttbar[bar/.style={fill=primaryCyan, draw=primaryCyan, rounded corners=1pt}, bar height=0.40, bar top shift=0.48, bar label font=\tiny\bfseries\color{white}]{F08}{2026-04-28}{2026-07-23} \\

    % F09 
    \ganttbar[bar/.style={fill=gray!45, draw=secondaryMagenta, line width=0.5pt, rounded corners=1pt}, bar height=0.40, bar top shift=0.08, bar label font=\tiny\bfseries\color{black!80}]{F09-Plan}{2026-04-15}{2026-05-17}
    \ganttbar[bar/.style={fill=secondaryMagenta, draw=secondaryMagenta, rounded corners=1pt}, bar height=0.40, bar top shift=0.48, bar label font=\tiny\bfseries\color{white}]{F09}{2026-05-04}{2026-10-17} \\

    % F10 
    \ganttbar[bar/.style={fill=gray!45, draw=secondaryMagenta, line width=0.5pt, rounded corners=1pt}, bar height=0.40, bar top shift=0.08, bar label font=\tiny\bfseries\color{black!80}]{F10-Plan}{2026-05-18}{2026-06-12}
    \ganttbar[bar/.style={fill=secondaryMagenta, draw=secondaryMagenta, rounded corners=1pt}, bar height=0.40, bar top shift=0.48, bar label font=\tiny\bfseries\color{white}]{F10}{2026-10-01}{2026-10-17}

    % Líneas semanales 
    \ganttvrule[vrule/.style={draw=gray!60, dotted, line width=0.4pt}, vrule label node/.style={anchor=south, yshift=7.22cm, font=\tiny\bfseries\color{gray!80!black}, fill=white, inner sep=0.5pt}]{19}{2026-01-19}
    \ganttvrule[vrule/.style={draw=gray!60, dotted, line width=0.4pt}, vrule label node/.style={anchor=south, yshift=7.22cm, font=\tiny\bfseries\color{gray!80!black}, fill=white, inner sep=0.5pt}]{26}{2026-01-26}
    \ganttvrule[vrule/.style={draw=gray!60, dotted, line width=0.4pt}, vrule label node/.style={anchor=south, yshift=7.22cm, font=\tiny\bfseries\color{gray!80!black}, fill=white, inner sep=0.5pt}]{02}{2026-02-02}
    \ganttvrule[vrule/.style={draw=gray!60, dotted, line width=0.4pt}, vrule label node/.style={anchor=south, yshift=7.22cm, font=\tiny\bfseries\color{gray!80!black}, fill=white, inner sep=0.5pt}]{09}{2026-02-09}
    \ganttvrule[vrule/.style={draw=gray!60, dotted, line width=0.4pt}, vrule label node/.style={anchor=south, yshift=7.22cm, font=\tiny\bfseries\color{gray!80!black}, fill=white, inner sep=0.5pt}]{16}{2026-02-16}
    \ganttvrule[vrule/.style={draw=gray!60, dotted, line width=0.4pt}, vrule label node/.style={anchor=south, yshift=7.22cm, font=\tiny\bfseries\color{gray!80!black}, fill=white, inner sep=0.5pt}]{23}{2026-02-23}
    \ganttvrule[vrule/.style={draw=gray!60, dotted, line width=0.4pt}, vrule label node/.style={anchor=south, yshift=7.22cm, font=\tiny\bfseries\color{gray!80!black}, fill=white, inner sep=0.5pt}]{02}{2026-03-02}
    \ganttvrule[vrule/.style={draw=gray!60, dotted, line width=0.4pt}, vrule label node/.style={anchor=south, yshift=7.22cm, font=\tiny\bfseries\color{gray!80!black}, fill=white, inner sep=0.5pt}]{09}{2026-03-09}
    \ganttvrule[vrule/.style={draw=gray!60, dotted, line width=0.4pt}, vrule label node/.style={anchor=south, yshift=7.22cm, font=\tiny\bfseries\color{gray!80!black}, fill=white, inner sep=0.5pt}]{16}{2026-03-16}
    \ganttvrule[vrule/.style={draw=gray!60, dotted, line width=0.4pt}, vrule label node/.style={anchor=south, yshift=7.22cm, font=\tiny\bfseries\color{gray!80!black}, fill=white, inner sep=0.5pt}]{23}{2026-03-23}
    \ganttvrule[vrule/.style={draw=gray!60, dotted, line width=0.4pt}, vrule label node/.style={anchor=south, yshift=7.22cm, font=\tiny\bfseries\color{gray!80!black}, fill=white, inner sep=0.5pt}]{30}{2026-03-30}
    \ganttvrule[vrule/.style={draw=gray!60, dotted, line width=0.4pt}, vrule label node/.style={anchor=south, yshift=7.22cm, font=\tiny\bfseries\color{gray!80!black}, fill=white, inner sep=0.5pt}]{06}{2026-04-06}
    \ganttvrule[vrule/.style={draw=gray!60, dotted, line width=0.4pt}, vrule label node/.style={anchor=south, yshift=7.22cm, font=\tiny\bfseries\color{gray!80!black}, fill=white, inner sep=0.5pt}]{13}{2026-04-13}
    \ganttvrule[vrule/.style={draw=gray!60, dotted, line width=0.4pt}, vrule label node/.style={anchor=south, yshift=7.22cm, font=\tiny\bfseries\color{gray!80!black}, fill=white, inner sep=0.5pt}]{20}{2026-04-20}
    \ganttvrule[vrule/.style={draw=gray!60, dotted, line width=0.4pt}, vrule label node/.style={anchor=south, yshift=7.22cm, font=\tiny\bfseries\color{gray!80!black}, fill=white, inner sep=0.5pt}]{27}{2026-04-27}
    \ganttvrule[vrule/.style={draw=gray!60, dotted, line width=0.4pt}, vrule label node/.style={anchor=south, yshift=7.22cm, font=\tiny\bfseries\color{gray!80!black}, fill=white, inner sep=0.5pt}]{04}{2026-05-04}
    \ganttvrule[vrule/.style={draw=gray!60, dotted, line width=0.4pt}, vrule label node/.style={anchor=south, yshift=7.22cm, font=\tiny\bfseries\color{gray!80!black}, fill=white, inner sep=0.5pt}]{11}{2026-05-11}
    \ganttvrule[vrule/.style={draw=gray!60, dotted, line width=0.4pt}, vrule label node/.style={anchor=south, yshift=7.22cm, font=\tiny\bfseries\color{gray!80!black}, fill=white, inner sep=0.5pt}]{18}{2026-05-18}
    \ganttvrule[vrule/.style={draw=gray!60, dotted, line width=0.4pt}, vrule label node/.style={anchor=south, yshift=7.22cm, font=\tiny\bfseries\color{gray!80!black}, fill=white, inner sep=0.5pt}]{25}{2026-05-25}
    \ganttvrule[vrule/.style={draw=gray!60, dotted, line width=0.4pt}, vrule label node/.style={anchor=south, yshift=7.22cm, font=\tiny\bfseries\color{gray!80!black}, fill=white, inner sep=0.5pt}]{01}{2026-06-01}
    \ganttvrule[vrule/.style={draw=gray!60, dotted, line width=0.4pt}, vrule label node/.style={anchor=south, yshift=7.22cm, font=\tiny\bfseries\color{gray!80!black}, fill=white, inner sep=0.5pt}]{08}{2026-06-08}
    \ganttvrule[vrule/.style={draw=gray!60, dotted, line width=0.4pt}, vrule label node/.style={anchor=south, yshift=7.22cm, font=\tiny\bfseries\color{gray!80!black}, fill=white, inner sep=0.5pt}]{15}{2026-06-15}
    \ganttvrule[vrule/.style={draw=gray!60, dotted, line width=0.4pt}, vrule label node/.style={anchor=south, yshift=7.22cm, font=\tiny\bfseries\color{gray!80!black}, fill=white, inner sep=0.5pt}]{22}{2026-06-22}
    \ganttvrule[vrule/.style={draw=gray!60, dotted, line width=0.4pt}, vrule label node/.style={anchor=south, yshift=7.22cm, font=\tiny\bfseries\color{gray!80!black}, fill=white, inner sep=0.5pt}]{29}{2026-06-29}
    \ganttvrule[vrule/.style={draw=gray!60, dotted, line width=0.4pt}, vrule label node/.style={anchor=south, yshift=7.22cm, font=\tiny\bfseries\color{gray!80!black}, fill=white, inner sep=0.5pt}]{06}{2026-07-06}
    \ganttvrule[vrule/.style={draw=gray!60, dotted, line width=0.4pt}, vrule label node/.style={anchor=south, yshift=7.22cm, font=\tiny\bfseries\color{gray!80!black}, fill=white, inner sep=0.5pt}]{13}{2026-07-13}
    \ganttvrule[vrule/.style={draw=gray!60, dotted, line width=0.4pt}, vrule label node/.style={anchor=south, yshift=7.22cm, font=\tiny\bfseries\color{gray!80!black}, fill=white, inner sep=0.5pt}]{20}{2026-07-20}
    \ganttvrule[vrule/.style={draw=gray!60, dotted, line width=0.4pt}, vrule label node/.style={anchor=south, yshift=7.22cm, font=\tiny\bfseries\color{gray!80!black}, fill=white, inner sep=0.5pt}]{27}{2026-07-27}
    \ganttvrule[vrule/.style={draw=gray!60, dotted, line width=0.4pt}, vrule label node/.style={anchor=south, yshift=7.22cm, font=\tiny\bfseries\color{gray!80!black}, fill=white, inner sep=0.5pt}]{03}{2026-08-03}
    \ganttvrule[vrule/.style={draw=gray!60, dotted, line width=0.4pt}, vrule label node/.style={anchor=south, yshift=7.22cm, font=\tiny\bfseries\color{gray!80!black}, fill=white, inner sep=0.5pt}]{10}{2026-08-10}
    \ganttvrule[vrule/.style={draw=gray!60, dotted, line width=0.4pt}, vrule label node/.style={anchor=south, yshift=7.22cm, font=\tiny\bfseries\color{gray!80!black}, fill=white, inner sep=0.5pt}]{17}{2026-08-17}
    \ganttvrule[vrule/.style={draw=gray!60, dotted, line width=0.4pt}, vrule label node/.style={anchor=south, yshift=7.22cm, font=\tiny\bfseries\color{gray!80!black}, fill=white, inner sep=0.5pt}]{24}{2026-08-24}
    \ganttvrule[vrule/.style={draw=gray!60, dotted, line width=0.4pt}, vrule label node/.style={anchor=south, yshift=7.22cm, font=\tiny\bfseries\color{gray!80!black}, fill=white, inner sep=0.5pt}]{31}{2026-08-31}
    \ganttvrule[vrule/.style={draw=gray!60, dotted, line width=0.4pt}, vrule label node/.style={anchor=south, yshift=7.22cm, font=\tiny\bfseries\color{gray!80!black}, fill=white, inner sep=0.5pt}]{07}{2026-09-07}
    \ganttvrule[vrule/.style={draw=gray!60, dotted, line width=0.4pt}, vrule label node/.style={anchor=south, yshift=7.22cm, font=\tiny\bfseries\color{gray!80!black}, fill=white, inner sep=0.5pt}]{14}{2026-09-14}
    \ganttvrule[vrule/.style={draw=gray!60, dotted, line width=0.4pt}, vrule label node/.style={anchor=south, yshift=7.22cm, font=\tiny\bfseries\color{gray!80!black}, fill=white, inner sep=0.5pt}]{21}{2026-09-21}
    \ganttvrule[vrule/.style={draw=gray!60, dotted, line width=0.4pt}, vrule label node/.style={anchor=south, yshift=7.22cm, font=\tiny\bfseries\color{gray!80!black}, fill=white, inner sep=0.5pt}]{28}{2026-09-28}
    \ganttvrule[vrule/.style={draw=gray!60, dotted, line width=0.4pt}, vrule label node/.style={anchor=south, yshift=7.22cm, font=\tiny\bfseries\color{gray!80!black}, fill=white, inner sep=0.5pt}]{05}{2026-10-05}
    \ganttvrule[vrule/.style={draw=gray!60, dotted, line width=0.4pt}, vrule label node/.style={anchor=south, yshift=7.22cm, font=\tiny\bfseries\color{gray!80!black}, fill=white, inner sep=0.5pt}]{12}{2026-10-12}
\end{ganttchart}%
}
\vspace{0.25cm}
\begin{center}
\footnotesize
\begin{tabular}{clccl}
\fcolorbox{black}{primaryMagenta}{\rule{0pt}{4pt}\rule{4pt}{0pt}} & Gestión y planificación & \hspace{0.5cm} &
\fcolorbox{black}{primaryCyan}{\rule{0pt}{4pt}\rule{4pt}{0pt}} & Desarrollo de código \\
\fcolorbox{black}{secondaryCyan}{\rule{0pt}{4pt}\rule{4pt}{0pt}} & Ciclo de vida y calidad & \hspace{0.5cm} &
\fcolorbox{black}{secondaryMagenta}{\rule{0pt}{4pt}\rule{4pt}{0pt}} & Documentación y cierre
\end{tabular}
\end{center}
\caption{Diagrama de Gantt de seguimiento: cronograma planificado vs. real.}
\label{fig:gantt_comparativa}
\end{figure}

\textcolor{red}{\textbf{[TODO: Establecer fin de fase 9 y comienzo de fase 10]}}

\section{Gastos incurridos frente a Línea Base de Costes}\label{sec:gastos_incurridos_costes}

Para evaluar las diferencias entre la \hyperref[sec:linea_base_de_costes]{Línea Base de Costes} y los gastos reales del proyecto, el análisis se centra exclusivamente en el CAPEX. El OPEX se mantiene idéntico a lo previsto, ya que los costes recurrentes de despliegue y consumo de servicios (incluida la API de IA para cuando entre en producción) se consideraron fijos para el proyecto y no han sufrido variaciones.

La variación del CAPEX se debe principalmente a dos factores derivados de las desviaciones temporales del proyecto: por un lado, el desvío en los costes laborales según el esfuerzo dedicado, y por otro, la extensión en el uso de las licencias de software. Al haberse trasladado el cierre del proyecto de junio a octubre, las suscripciones de GitHub Team y Microsoft Teams Essentials debieron mantenerse activas durante cuatro meses adicionales (pasando de 6 a 10 meses), mientras que la licencia comercial anual de Google AI Premium y el equipamiento no experimentaron cambios.

En la siguiente tabla se recoge el balance consolidado del presupuesto de capital frente a los costes reales:

\begin{table}[H]
\centering
\small
\renewcommand{\arraystretch}{1.2}
\makebox[\textwidth][c]{%
\begin{tabular}{lrrr}
\hline
\textbf{Concepto} & \textbf{Coste previsto} & \textbf{Coste real} & \textbf{Desviación} \\ \hline
Costes laborales                                        & 11.552,65 \euro & [   \euro] & [   \euro] \\
Equipamiento informático                                & 500,00 \euro    & 500,00 \euro & 0,00 \euro \\
Licencia Google AI Premium (anual)                      & 218,07 \euro    & 218,07 \euro & 0,00 \euro \\
Licencia GitHub Team                                    & 21,82 \euro{} (6 m)   & 36,37 \euro{} (10 m) & 14,55 \euro \\
Licencia Microsoft Teams Essentials                     & 21,82 \euro{} (6 m)   & 36,37 \euro{} (10 m) & 14,55 \euro \\
Reserva de contingencia                                 & 12.314,36 \euro & [   \euro] & [   \euro] \\ \hline
\textbf{Total CAPEX}                                    & \textbf{13.545,80 \euro} & \textbf{[   \euro]} & \textbf{[   \euro]} \\ \hline
\end{tabular}%
}
\caption{Comparativa de Gastos de Capital (CAPEX) presupuestados vs. reales.}
\label{tab:resumen_capex_real}
\end{table}

\textcolor{red}{\textbf{[TODO:Completar tabla]}}

Para comprender de dónde procede la desviación observada en la partida de costes laborales, es necesario descender al detalle del esfuerzo por rol profesional. La redistribución de horas entre tareas alteró la dedicación asignada a cada perfil. A continuación se desglosan las horas e importes correspondientes a partir de las tarifas por hora fijadas en la planificación inicial:

\begin{table}[H]
\centering
\small
\renewcommand{\arraystretch}{1.3}
\noindent\makebox[\textwidth][c]{%
\begin{tabular}{lccccc}
\hline
\textbf{Rol} & \textbf{Precio/hora} & \textbf{Horas previstas} & \textbf{Horas reales} & \textbf{Coste previsto} & \textbf{Coste real} \\ \hline
Project Manager      & 46,09 \euro & 18 h  & [   h] & 829,62 \euro   & [   \euro] \\
Analista             & 38,33 \euro & 87 h  & [   h] & 3.334,71 \euro & [   \euro] \\
Diseñador UI         & 33,43 \euro & 22 h  & [   h] & 735,46 \euro   & [   \euro] \\
Programador Backend  & 23,85 \euro & 79 h  & [   h] & 1.884,15 \euro & [   \euro] \\
Programador Frontend & 23,85 \euro & 47 h  & [   h] & 1.120,95 \euro & [   \euro] \\
Tester               & 23,68 \euro & 15 h  & [   h] & 355,20 \euro   & [   \euro] \\
Técnico              & 19,58 \euro & 32 h  & [   h] & 626,56 \euro   & [   \euro] \\ \hline
\textbf{Total}       & --          & \textbf{300 h} & \textbf{[   h]} & \textbf{8.886,65 \euro} & \textbf{[   \euro]} \\ \hline
\end{tabular}%
}
\caption{Comparativa de dedicación y costes de personal por rol profesional.}
\label{tab:horas_roles_previstas_reales}
\end{table}

\textcolor{red}{\textbf{[TODO: Completar tabla y redactar pequeña retrospectiva final de desviaciones de los roles]}}

Como consecuencia de esta variación en el gasto de capital y al mantenerse inmutable el OPEX, el coste total de propiedad (TCO) del proyecto pasa de los 15.198,56 € estimados inicialmente a un valor definitivo de [   \euro], cerrando la auditoría económica del desarrollo.

\textcolor{red}{\textbf{[TODO: Añadir tco final]}}

\section{Desviaciones y gestión de cambios}\label{sec:desviaciones_y_gestion_cambios}

A lo largo del ciclo de vida del proyecto surgieron contingencias técnicas, restricciones de tiempo y oportunidades de mejora que hicieron necesario activar mecanismos formales de control del cambio. Con el objetivo de proteger la calidad técnica del producto software y evitar desviaciones incontroladas de la planificación, cada alteración significativa respecto a la línea base se analizó y documentó mediante una ficha formal de solicitud de cambio (\textit{Change Request}), evaluando su impacto directo sobre el alcance, el cronograma y los costes.

\subsection{Registro de cambios}\label{sub:registro_cambios}

En este apartado se recoge el historial de las principales decisiones de cambio gestionadas durante el desarrollo, estructuradas como \textit{Change Requests} (CR). Cada ficha documenta el estado de la solicitud, el contexto que la motivó, la decisión adoptada y las consecuencias derivadas de su resolución.

\begin{table}[H]
\centering
\renewcommand{\arraystretch}{1.5}
\setlength{\tabcolsep}{12pt}
\makebox[\textwidth][c]{%
\begin{tabular}{lp{11.5cm}}
\hline
\textbf{CR} & CR-01: Aplazamiento de la defensa a convocatoria de octubre. \\
\hline
\textbf{Estado} & Aprobado. \\
\hline
\textbf{Contexto} & El esfuerzo necesario para estabilizar las pruebas de integración en WebSockets y redactar la memoria con el rigor exigido impidió concluir en junio sin comprometer la calidad. \\
\hline
\textbf{Decisión} & Trasladar formalmente la entrega de la memoria y la defensa técnica de la convocatoria ordinaria de junio a la extraordinaria de octubre de 2026. \\
\hline
\textbf{Impacto en alcance} & Ninguno. \\
\hline
\textbf{Impacto en tiempo} & Ampliación de 4 meses en la ventana temporal global del proyecto. \\
\hline
\textbf{Impacto en coste} & Incremento de 29,10 \euro{} en licencias de desarrollo CAPEX por 4 meses adicionales de GitHub Team y Microsoft Teams Essentials. \\
\hline
\end{tabular}%
}
\caption{Registro de cambios - CR-01.}
\label{tab:cr_01_aplazamiento}
\end{table}

\begin{table}[H]
\centering
\renewcommand{\arraystretch}{1.5}
\setlength{\tabcolsep}{12pt}
\makebox[\textwidth][c]{%
\begin{tabular}{lp{11.5cm}}
\hline
\textbf{CR} & CR-02: Aplazamiento de funcionalidades del asistente de generación con IA (\hyperref[obj:09]{OBJ-09}). \\
\hline
\textbf{Estado} & Aprobado. \\
\hline
\textbf{Contexto} & La latencia y los límites de consumo (\textit{rate limits}) de la API de Gemini planteaban riesgos de bloqueo operativo en vivo, exigiendo un esfuerzo desproporcionado para asegurar un comportamiento determinista. \\
\hline
\textbf{Decisión} & Desacoplar la generación con IA del alcance actual, registrarla como línea de trabajo futuro en GitHub y priorizar la creación e importación manual de cuestionarios. \\
\hline
\textbf{Impacto en alcance} & Reducción del alcance funcional inmediato, preservando íntegro el flujo estándar de gestión de cuestionarios. \\
\hline
\textbf{Impacto en tiempo} & Liberación de horas de desarrollo backend y frontend, mitigando retrasos en el núcleo del sistema. \\
\hline
\textbf{Impacto en coste} & Derivado de la redistribución de horas entre roles profesionales. \\
\hline
\end{tabular}%
}
\caption{Registro de cambios - CR-02.}
\label{tab:cr_02_ia}
\end{table}

\begin{table}[H]
\centering
\renewcommand{\arraystretch}{1.5}
\setlength{\tabcolsep}{12pt}
\makebox[\textwidth][c]{%
\begin{tabular}{lp{11.5cm}}
\hline
\textbf{CR} & CR-03: Postergación de la creación de grupos de alumnos (\hyperref[obj:11]{OBJ-11}). \\
\hline
\textbf{Estado} & Aprobado. \\
\hline
\textbf{Contexto} & Se priorizó centrar el tiempo de desarrollo en la estabilidad de las partidas en tiempo real antes que en la gestión de grupos. \\
\hline
\textbf{Decisión} & Posponer la gestión de grupos de alumnos y dejarla registrada en GitHub como trabajo futuro. \\
\hline
\textbf{Impacto en alcance} & Se excluye la funcionalidad de grupos de alumnos en esta entrega para asegurar el funcionamiento del juego en vivo. \\
\hline
\textbf{Impacto en tiempo} & Se liberan horas de desarrollo para dedicarlas a la estabilidad y pruebas del sistema. \\
\hline
\textbf{Impacto en coste} & Derivado de la redistribución de horas entre roles profesionales. \\
\hline
\end{tabular}%
}
\caption{Registro de cambios - CR-03.}
\label{tab:cr_03_obj11}
\end{table}

\begin{table}[H]
\centering
\renewcommand{\arraystretch}{1.5}
\setlength{\tabcolsep}{12pt}
\makebox[\textwidth][c]{%
\begin{tabular}{lp{11.5cm}}
\hline
\textbf{CR} & CR-04: Ampliación de la batería y estrategia de pruebas. \\
\hline
\textbf{Estado} & Aprobado. \\
\hline
\textbf{Contexto} & La estrategia de pruebas planificada inicialmente resultó insuficiente para asegurar la calidad global del sistema, haciendo necesario un enfoque de verificación más exhaustivo. \\
\hline
\textbf{Decisión} & Rediseñar y ampliar la estrategia de pruebas, incorporando una batería más completa de casos de prueba tanto a nivel unitario como de integración. \\
\hline
\textbf{Impacto en alcance} & Incremento sustancial en la cobertura de pruebas y en la robustez general de la aplicación. \\
\hline
\textbf{Impacto en tiempo} & Aumento notable de horas dedicadas a la fase de pruebas frente a las 15 horas presupuestadas originalmente. \\
\hline
\textbf{Impacto en coste} & Incremento en el coste imputado al rol de Tester dentro del balance contable del CAPEX. \\
\hline
\end{tabular}%
}
\caption{Registro de cambios - CR-04.}
\label{tab:cr_04_testing}
\end{table}

\begin{table}[H]
\centering
\renewcommand{\arraystretch}{1.5}
\setlength{\tabcolsep}{12pt}
\makebox[\textwidth][c]{%
\begin{tabular}{lp{11.5cm}}
\hline
\textbf{CR} & CR-05: Incorporación de Redis para la sincronización de la sala. \\
\hline
\textbf{Estado} & Desestimado. \\
\hline
\textbf{Contexto} & Se evaluó utilizar Redis para guardar el estado de las salas y permitir que varios servidores pudieran comunicarse entre sí ante un gran número de usuarios simultáneos. \\
\hline
\textbf{Decisión} & Desestimar el uso de Redis y mantener la gestión de las salas en la propia memoria del servidor actual. \\
\hline
\textbf{Impacto en alcance} & Ninguno; se mantiene una arquitectura más simple y directa, suficiente para las necesidades del proyecto. \\
\hline
\textbf{Impacto en tiempo} & Se evita un trabajo extra estimado entre 10 y 15 horas de desarrollo y configuración. \\
\hline
\textbf{Impacto en coste} & Ninguno; además, se evita el coste mensual extra de contratar un servicio adicional en la nube. \\
\hline
\end{tabular}%
}
\caption{Registro de cambios - CR-05.}
\label{tab:cr_05_redis}
\end{table}

\begin{table}[H]
\centering
\renewcommand{\arraystretch}{1.5}
\setlength{\tabcolsep}{12pt}
\makebox[\textwidth][c]{%
\begin{tabular}{lp{11.5cm}}
\hline
\textbf{CR} & CR-06: Despliegue efectivo de SonarCloud y refuerzo de controles en local (\textit{pre-commits}). \\
\hline
\textbf{Estado} & Aprobado. \\
\hline
\textbf{Contexto} & La activación en CI/CD del análisis con SonarCloud se demoró; al procesar el código acumulado, el volumen de avisos obligó a atajar las incidencias antes de subirlas al repositorio. \\
\hline
\textbf{Decisión} & Habilitar formalmente el análisis continuo en GitHub Actions e implantar ganchos de \textit{pre-commit} junto a formateadores locales para filtrar defectos de estilo y tipado en local. \\
\hline
\textbf{Impacto en alcance} & Incorporación de controles estáticos automatizados en el entorno local y en la integración continua. \\
\hline
\textbf{Impacto en tiempo} & Dedicación técnica imprevista en configurar las herramientas locales y solventar los avisos acumulados. \\
\hline
\textbf{Impacto en coste} & Ninguno en licencias (uso del plan gratuito de SonarCloud); imputación adicional de horas al rol Técnico. \\
\hline
\end{tabular}%
}
\caption{Registro de cambios - CR-06.}
\label{tab:cr_06_sonar}
\end{table}

\begin{table}[H]
\centering
\renewcommand{\arraystretch}{1.5}
\setlength{\tabcolsep}{12pt}
\makebox[\textwidth][c]{%
\begin{tabular}{lp{11.5cm}}
\hline
\textbf{CR} & CR-07: Reducción de deuda técnica en SonarCloud y congelación del alcance. \\
\hline
\textbf{Estado} & Aprobado. \\
\hline
\textbf{Contexto} & La integración tardía del análisis de SonarCloud generó un volumen acumulado de deuda técnica que exigió atención inmediata, poniendo en riesgo la fecha límite de entrega si se continuaba desarrollando nuevas funcionalidades. \\
\hline
\textbf{Decisión} & Congelar el desarrollo de nuevas características para priorizar la corrección de errores críticos en SonarCloud, posponiendo los requisitos de mejora no esenciales como tareas de trabajo futuro en GitHub, en consonancia con la retrospectiva del alcance. \\
\hline
\textbf{Impacto en alcance} & Se detiene la incorporación de requisitos de mejora secundarios, asegurando una base de código limpia y libre de fallos graves para el núcleo de la plataforma. \\
\hline
\textbf{Impacto en tiempo} & Permite cumplir con los plazos de entrega y dedicar el tiempo necesario a la redacción rigurosa de la memoria técnica. \\
\hline
\textbf{Impacto en coste} & Derivado de la redistribución de horas entre roles profesionales. \\
\hline
\end{tabular}%
}
\caption{Registro de cambios - CR-07.}
\label{tab:cr_07_backlog}
\end{table}

Este conjunto de solicitudes de cambio pone de manifiesto que las desviaciones temporales y económicas registradas no fueron fruto de la improvisación, sino de decisiones conscientes de ingeniería orientadas a priorizar la estabilidad técnica, la verificación exhaustiva y la calidad del producto final frente a las restricciones temporales.

\section{Resultados de las pruebas y validación}\label{sec:resultados_pruebas}

Una vez definida la estrategia de verificación en la fase de planificación y detallada la infraestructura técnica de ejecución en el capítulo de implementación, en este apartado se exponen exclusivamente los resultados obtenidos. El objetivo es contrastar las métricas reales de cobertura, robustez y rendimiento frente a los umbrales de aceptación fijados como criterio de calidad para el proyecto. Para facilitar la trazabilidad completa, el material probatorio complementario (reportes interactivos, registros íntegros y capturas) se encuentra recopilado en el anexo de \hyperref[anexo:registro_pruebas_cobertura]{registro de pruebas de cobertura}.

\subsection{Pruebas unitarias y de integración del backend}\label{sub:resultados_pruebas_backend}

La suite del \textit{backend} cubre tanto la lógica de negocio aislada (modelos, servicios y esquemas) como la integración de los puntos de entrada de la API contra la base de datos y el motor de autenticación. 

\begin{evidence}[H]
\begin{lstlisting}
tests\backend\test_integration_content.py (*@\covGreen{....}@*)                  (*@\covGreen{[\ \ 5\%]}@*)
tests\backend\test_integration_stage.py (*@\covGreen{...........}@*)             (*@\covGreen{[\ 21\%]}@*)
tests\backend\test_integration_users.py (*@\covGreen{....}@*)                    (*@\covGreen{[\ 26\%]}@*)
tests\backend\test_unit_content.py (*@\covGreen{.............}@*)                (*@\covGreen{[\ 45\%]}@*)
tests\backend\test_unit_stage.py (*@\covGreen{......................}@*)         (*@\covGreen{[\ 76\%]}@*)
tests\backend\test_unit_users.py (*@\covGreen{.................}@*)              (*@\covGreen{[100\%]}@*)

Name                                             Stmts   Miss   Cover
-------------------------------------------------------------------------
backend\routers\content.py                         188      0  (*@\covGreen{100\%}@*)
backend\routers\stage.py                           454      9   (*@\covGreen{98\%}@*)
backend\routers\users.py                           211      0  (*@\covGreen{100\%}@*)
Otros (auth, database, models, schemas...)         398     23   (*@\covGreen{94\%}@*)
-------------------------------------------------------------------------
TOTAL                                             1251     32   (*@\covGreen{97\%}@*)

(*@\covGreen{========================== 71 passed in 13.50s ==========================}@*)
\end{lstlisting}
\caption{Resumen de ejecución y cobertura en backend.}
\label{ev:resumen_backend}
\end{evidence}

La batería completó la ejecución de 71 pruebas sin registrar ningún fallo en un tiempo de 13{,}50 segundos. A nivel de cobertura:

\begin{itemize}
    \item Se alcanzó una cobertura agregada de sentencias del \textbf{97\,\%}, superando ampliamente el umbral mínimo del 80\,\% fijado en los \hyperref[tab:umbrales_cobertura]{umbrales de cobertura}.
    \item Los controladores de contenidos (\texttt{content.py}) y de usuarios (\texttt{users.py}) alcanzaron una cobertura plena del \textbf{100\,\%}.
    \item La máquina de estados de las salas (\texttt{stage.py}) registró un \textbf{98\,\%}, garantizando la solidez de las dinámicas en tiempo real frente a condiciones imprevistas.
\end{itemize}

Estos resultados quedan respaldados en el anexo mediante la captura de la \hyperref[ev:terminal_test_backend]{salida de terminal de las pruebas de \textbf{pytest}} y el informe de \hyperref[ev:cobertura_backend]{cobertura del \textit{backend}} generado en HTML con el desglose detallado por archivos.

\subsection{Pruebas de componentes del frontend}\label{sub:resultados_pruebas_frontend}

En la interfaz visual, las pruebas de componentes permiten verificar el renderizado, el ciclo de vida, la gestión de estados reactivos y la emisión de eventos de forma aislada.

\begin{evidence}[H]
\begin{lstlisting}
Test Files  (*@\covGreen{25 passed}@*) (25)
     Tests  (*@\covGreen{193 passed}@*) (193)

-----------------------|---------|----------|---------|---------|
File                   | % Stmts | % Branch | % Funcs | % Lines |
-----------------------|---------|----------|---------|---------|
All files              |   (*@\covGreen{90.13}@*) |    (*@\covGreen{80.29}@*) |   (*@\covGreen{88.56}@*) |   (*@\covGreen{92.68}@*) |
 auth                  |   (*@\covGreen{93.37}@*) |    (*@\covGreen{84.93}@*) |   (*@\covGreen{95.16}@*) |   (*@\covGreen{94.17}@*) |
 management            |   (*@\covGreen{87.82}@*) |    (*@\covGreen{76.37}@*) |   (*@\covGreen{86.36}@*) |   (*@\covGreen{90.78}@*) |
 room-access           |   (*@\covGreen{94.91}@*) |     (*@\covGreen{75.8}@*) |   (*@\covGreen{94.73}@*) |   (*@\covGreen{97.36}@*) |
 room-play             |   (*@\covGreen{89.66}@*) |    (*@\covGreen{82.47}@*) |   (*@\covGreen{86.9}@*)  |   (*@\covGreen{92.83}@*) |
-----------------------|---------|----------|---------|---------|
\end{lstlisting}
\caption{Resumen de pruebas y cobertura de componentes.}
\label{ev:resumen_componentes}
\end{evidence}

La ejecución completó los 193 casos de prueba a lo largo de los 25 ficheros evaluados. El análisis de las métricas arroja las siguientes conclusiones sobre la robustez visual:

\begin{itemize}
    \item El umbral fijado del 80\,\% se superó de forma generalizada, alcanzando un \textbf{90{,}13\,\%} en sentencias y consolidando un \textbf{92{,}68\,\%} sobre el total de líneas de código fuente.
    \item La cobertura de ramas se situó en el \textbf{80{,}29\,\%}, certificando el correcto tratamiento de los flujos alternativos del cliente, tales como errores de red, validaciones de entrada o pantallas intermedias de espera.
    \item Los dominios de acceso e incorporación (\texttt{auth} y \texttt{room-access}) registraron los valores más altos (superando el 94\,\% de cobertura), blindando los puntos iniciales de entrada de la experiencia de usuario.
\end{itemize}

El detalle completo de esta batería se puede contrastar en el anexo consultando la \hyperref[ev:terminal_test_componentes]{salida por consola de \textbf{Vitest}} junto con la perspectiva del \hyperref[ev:cobertura_componentes]{reporte HTML de cobertura de componentes}.

\subsection{Pruebas de extremo a extremo (E2E)}\label{sub:resultados_pruebas_e2e}

Las pruebas de extremo a extremo verifican los flujos completos de usuario sobre instancias reales de navegador, asegurando la correcta integración entre la interfaz, los servicios de la API y los canales de comunicación bidireccional. 
\begin{evidence}[H]
\begin{lstlisting}
Test Files  (*@\covGreen{10 passed}@*) (10)
     Tests  (*@\covGreen{112 passed}@*) (112)

-------------------|---------|---------|----------|---------|---------|
Name               | % Bytes | % Stms | % Branch | % Funcs | % Lines |
-------------------|---------|---------|----------|---------|---------|
auth               |  (*@\covGreen{95.44}@*)  |   (*@\covGreen{95.53}@*) |    (*@\covGreen{91.77}@*) |   (*@\covGreen{85.25}@*) |   (*@\covGreen{94.79}@*) |
management         |  (*@\covGreen{93.58}@*)  |   (*@\covGreen{88.37}@*) |    (*@\covGreen{80.18}@*) |   (*@\covGreen{92.35}@*) |   (*@\covGreen{92.37}@*) |
room-access        |  (*@\covGreen{97.80}@*)  |   (*@\covGreen{90.00}@*) |    (*@\covGreen{83.05}@*) |  (*@\covGreen{100.00}@*) |   (*@\covGreen{95.22}@*) |
room-play          |  (*@\covGreen{92.04}@*)  |   (*@\covGreen{91.88}@*) |    (*@\covGreen{84.10}@*) |   (*@\covGreen{92.11}@*) |   (*@\covGreen{90.65}@*) |
-------------------|---------|---------|----------|---------|---------|
Summary            |  (*@\covGreen{93.81}@*)  |   (*@\covGreen{91.30}@*) |    (*@\covGreen{83.92}@*) |   (*@\covGreen{91.41}@*) |   (*@\covGreen{92.58}@*) |
-------------------|---------|---------|----------|---------|---------|
\end{lstlisting}
\caption{Resumen de pruebas y cobertura \textit{End-to-End}.}
\label{ev:resumen_e2e}
\end{evidence}

La ejecución completó todos los escenarios sin registrar fallos sobre instancias reales de navegador. Más allá del volumen de pruebas aprobadas, la relevancia de este nivel reside en verificar la interoperabilidad entre la interfaz, los servicios de la API y los canales de comunicación en tiempo real. En este sentido, la cobertura global del 91{,}30\,\% en sentencias y del 83{,}92\,\% en \textit{branches} valida la consistencia de los flujos críticos de la plataforma, desde la persistencia de sesión en \texttt{auth} hasta la sincronización reactiva durante la partida en \texttt{room-play}.

Destaca el comportamiento del módulo \texttt{room-access}, que alcanzó el 100\,\% en cobertura de funciones; esto garantiza que todos los métodos asociados a la validación de salas, control de estados y admisión de participantes se ejercitan tal y como lo experimenta el usuario final. La traza detallada de todos los recorridos y el análisis visual de cobertura se recogen en el anexo mediante el \hyperref[ev:reporte_playwright]{reporte de ejecución de \textbf{Playwright}} y el \hyperref[ev:cobertura_e2e]{informe de cobertura de \textbf{Monocart}}.

\subsection{Pruebas de carga y rendimiento}\label{sub:resultados_pruebas_carga}

Las pruebas de rendimiento evalúan el comportamiento y los tiempos de respuesta del sistema bajo condiciones de concurrencia simulada. La batería se ejecutó de manera automatizada en modo desatendido, procesando las series temporales exportadas para contrastar las métricas obtenidas con los criterios de aceptación fijados.

\begin{evidence}[H]
\begin{lstlisting}
=======================================================================
 === INFORME DE EVALUACION DE UMBRALES DE RENDIMIENTO (P90) ===
=======================================================================
 Metrica / Endpoint                  | Valor P90 | Umbral    | Estado
-----------------------------------------------------------------------
 GET /content/quizzes/                | 20 ms   | <= 500 ms   | (*@\covGreen{[OK]}@*)
 POST /content/quizzes/               | 38 ms   | <= 500 ms   | (*@\covGreen{[OK]}@*)
 GET /content/quizzes/{id}            | 17 ms   | <= 200 ms   | (*@\covGreen{[OK]}@*)
 PUT /content/quizzes/{id}            | 22 ms   | <= 200 ms   | (*@\covGreen{[OK]}@*)
 POST /stage/participants             | 38 ms   | <= 200 ms   | (*@\covGreen{[OK]}@*)
 GET /stage/rooms                     | 69 ms   | <= 200 ms   | (*@\covGreen{[OK]}@*)
 POST /stage/rooms/                   | 84 ms   | <= 200 ms   | (*@\covGreen{[OK]}@*)
 GET /stage/rooms/verify/{join_code}  | 49 ms   | <= 200 ms   | (*@\covGreen{[OK]}@*)
 GET /stage/rooms/{id}                | 35 ms   | <= 200 ms   | (*@\covGreen{[OK]}@*)
 GET /stage/rooms/{id}/participants   | 12 ms   | <= 200 ms   | (*@\covGreen{[OK]}@*)
 POST /users/login                    | 240 ms  | <= 500 ms   | (*@\covGreen{[OK]}@*)
 GET /users/me                        | 12 ms   | <= 200 ms   | (*@\covGreen{[OK]}@*)
 WebSocket Connect                    | 28 ms   | <= 200 ms   | (*@\covGreen{[OK]}@*)
 Tasa de Errores (Global)             | 0.00%  | < 1.0%      | (*@\covGreen{[OK]}@*)
-----------------------------------------------------------------------
RESULTADO: TODAS LAS PRUEBAS DE RENDIMIENTO SUPERARON LOS UMBRALES (*@\covGreen{[OK]}@*)
=======================================================================
\end{lstlisting}
\caption{Informe de evaluación de umbrales de latencia $P_{90}$ y tasa de error.}
\label{ev:resumen_umbrales}
\end{evidence}

A lo largo de la simulación no se registró ningún fallo en las peticiones ni interrupciones en los canales de comunicación, consolidando una tasa de error del 0{,}00\,\% sobre el total de operaciones. En términos de latencia, todos los puntos de entrada cumplieron holgadamente con sus umbrales de aceptación: los endpoints de consulta y sincronización mantuvieron tiempos de respuesta $P_{90}$ inferiores a los 50\,ms, e incluso operaciones de mayor coste computacional como la autenticación (\texttt{/users/login}, con 240\,ms) se situaron muy por debajo del límite estricto de 500\,ms fijado en los \hyperref[tab:metricas_rendimiento]{umbrales de las pruebas de rendimiento}. Destaca el comportamiento del canal bidireccional (\textit{WebSocket Connect}), cuyo $P_{90}$ se consolidó en apenas 28\,ms bajo ráfagas concurrentes de conexión.

La trazabilidad completa del procedimiento se respalda en el anexo: la ejecución de la batería queda reflejada en la \hyperref[ev:terminal_test_locust]{captura de terminal de las pruebas de carga}, cuyos resultados brutos se recopilaron en los registros CSV correspondientes y fueron evaluados posteriormente mediante la \hyperref[ev:terminal_umbrales_locust]{ejecución del script de verificación de umbrales} de rendimiento.

\section{Análisis estático de código (SonarQube)}

% \section{Evaluación global frente a umbrales de calidad}\label{sub:evaluacion_umbrales_calidad}

\section{Retrospectiva técnica}
Reflexión sobre la implementación de las prácticas de CI/CD, testing y métricas de calidad, y decir si he cumplido los objetivos como las métricas de cobertura o de issues de sonar
